\documentclass[conference]{IEEEtran}
\usepackage[T1]{fontenc}
\IEEEoverridecommandlockouts
% The preceding line is only needed to identify funding in the first footnote. If that is unneeded, please comment it out.
\usepackage{cite}
\usepackage{amsmath,amssymb,amsfonts}
\usepackage{algorithmic}
\usepackage{graphicx}
\usepackage{textcomp}
\usepackage{xcolor}
\usepackage{booktabs}
\usepackage{amssymb}
\def\BibTeX{{\rm B\kern-.05em{\sc i\kern-.025em b}\kern-.08em
    T\kern-.1667em\lower.7ex\hbox{E}\kern-.125emX}}
\begin{document}

\title{ApproxiMPC: Distilling MPCs to LSTMs in F1Tenth}

\author{\IEEEauthorblockN{ Devin DeCosmo}
\IEEEauthorblockA{\textit{Mechanical Engineering} \\
\textit{Carnegie Mellon University}\\
Pittsburgh, PA \\
ddecosmo@andrew.cmu}
\and
\IEEEauthorblockN{ Regine DeCossard}
\IEEEauthorblockA{\textit{Computer Science} \\
\textit{Carnegie Mellon University}\\
Pittsburgh, PA \\
rdecossa@andrew.cmu}
\and
\IEEEauthorblockN{ Henry Feldhaus}
\IEEEauthorblockA{\textit{Mechanical Engineering} \\
\textit{Carnegie Mellon University}\\
Pittsburgh, PA \\
hfeldhau@andrew.cmu}
\and
\IEEEauthorblockN{ Ria Hegde}
\IEEEauthorblockA{\textit{Electrical Engineering} \\
\textit{Carnegie Mellon University}\\
Pittsburgh, PA \\
rhegde@andrew.cmu}
}

\maketitle

\begin{abstract}
UPDATE ME! Model Predictive Control (MPC) is widely used in autonomous racing because it produces high-performance control while enforcing system constraints. However, solving an optimization problem at every timestep introduces significant computational overhead, making it difficult to deploy on resource-constrained platforms such as F1Tenth vehicles. This work explores distilling an MPC controller into a Long Short-Term Memory (LSTM) network to reduce runtime cost while maintaining performance. Using a teacher-student approach, data generated from an MPC controller is used to train an LSTM to map sequential observations to control actions. The learned policy is evaluated in simulation against the MPC baseline using metrics such as lap time, collision frequency, steering smoothness, and inference time. The goal is to determine whether an LSTM can approximate MPC behavior while enabling faster, real-time control suitable for embedded systems.
\end{abstract}

\begin{IEEEkeywords}
component, formatting, style, styling, insert
\end{IEEEkeywords}

\section{Introduction - Condense}
Autonomous racing requires both fast decision-making and accurate control, while also being constrained by limited onboard computation. In platforms such as F1Tenth, this tradeoff becomes especially apparent, since control policies must run in real time on embedded hardware.

Model Predictive Control (MPC) is a strong approach in this setting because it optimizes control inputs while explicitly enforcing system constraints such as vehicle dynamics and friction limits. This typically leads to high-performance behavior, particularly for objectives like minimizing lap time. However, MPC requires solving an optimization problem at every timestep, which introduces significant computational overhead and limits its practicality for real-time deployment.

Long Short-Term Memory (LSTM) networks provide an alternative by approximating sequential control policies with low inference cost. Because they operate on temporal data, they can model system behavior and generate control outputs efficiently. The main limitation is that they do not explicitly enforce constraints, and their performance depends heavily on the quality of training data.

This work explores distilling an MPC controller into an LSTM policy within the F1Tenth environment. The goal is to retain the performance of MPC while reducing runtime computation, enabling a more efficient control policy for real-time autonomous racing.

\section{Related Work}

\subsection{Model Predictive Control}

There are a large number of control policies possible within F1Tenth but the most desirable of these are able to optimize a vehicles performance for minimum lap time or overall performance. 

Model Predictive Control (MPC) is a commonly used control policy within F1tenth and beyond, as it enables optimal or near optimal control in real time by optimizing an objective function while considering the capabilities of the system \cite{evans2024unifying}. MPC does this by combining input states and constraints to optimize an objective function, in this case minimum lap time. Input variables are split between state and control parameters. State inputs include the \textit{x} and \textit{y} position of the vehicle along with its heading angle, $\theta$. The main control inputs are velocity \textit{v}, steering angle $\delta$, and the center line speed $\dot{s}$ \cite{evans2024unifying}. Constraints are applied to the friction between the wheels and the surface, and steering regularization is applied. These techniques ensure the vehicle remains controllable and steering is minimized, allowing for a smoother final solution. The objective function can then be optimized as the car runs in both simulation and reality, maximizing vehicle performance. 

 MPC has been applied for several different use cases within F1Tenth, including general performance and lap time reduction, overtaking and opponent following, and variable friction conditions \cite{cataffo2022nonlinear,bhargav2021track,li2023realtime,nagy2023adaptive}. These MPC applications for F1Tenth use the dynamic bicycle model, allowing for high performance and accuracy with decreased degrees of freedom, DOFs \cite{cataffo2022nonlinear,bhargav2021track,li2023realtime,nagy2023adaptive}.

While MPC is a highly capable control policy, it requires a large amount of computation at run time to calculate the objective including the input states and constraints. This means that applications, such as F1Tenth, that are limited by computational power need to find workarounds to get performance from MPC. Strategies to optimize MPC for local performance like offline optimization reduce the reactivity of the control policy which can cause issues at high speeds or with dynamic obstacles like opponent vehicles. 

\subsection{Long Short Term Memory Models}

Long Short Term Memory Models, LSTMs, are a subset of Recurrent Neural Network models, RNNs. Like other RNNs, LSTMs take inputs as a sequence allowing them to consider temporal inputs and make predictions based on both past and present information \cite{ai6090215}. %Figure \ref{fig:LSTM} shows the overall architecture of an LSTM block. 

An important architectural part of an LSTM is the presence of forget, update, and output gates allowing for inputs to be considered only if they are relevant or useful making them adaptable to noise or errors in input data \cite{ai6090215}.

% \begin{figure}[h] % [h] forces the figure to be placed "here"
%     \centering
%     \includegraphics[width=0.3\textwidth]{figures/LSTM overview.png}
%     \caption{LSTM Architecture Overview}
%     \label{fig:LSTM} % Used for cross-referencing in the text
% \end{figure}

The modular nature of LSTM blocks also means that a model can be adapted to be lightweight enough to run on edge hardware like Nvidia Jetson Nanos. 

LSTMs have been tested using Formula Student races as training data and found they are robust, but require large amounts of training data \cite{srinivasan2020end}. LSTMs have also been tested on the F1tenth system using public datasets for velocity estimation, and MPC data for overtaking in simulation \cite{wegrzynowski2024learning,zhang2023f1tenth}. 

These results show that LSTMs are trainable and useful for creating control policies within F1tenth. 

\subsection{Transfer Learning MPC to LSTM}

While past studies have used MPC, LSTMs, and even trained LSTMs for specific tasks using data from MPC, none of these have focused on overall racing performance or been tested on a physical F1tenth vehicle. 

Additional work has discussed how MPC data can be curated and used to train LSTMs to mimic MPC performance \cite{munoz2023deep}. Additional work with the F1Tenth system focused on purely overtaking using LSTMs has also shown that MPC data can be collected and used to train LSTMs \cite{zhang2023f1tenth}.

These past works demonstrate that an existing MPC control policy can be used to collect racing data and train an LSTM to approach optimal performance. The moderate parameter size of LSTMs also ensure that training can be done efficiently, while still generalizing across race situations. 

\section{Experiment and Results}

- see bullets for (very) high level steering notes

\subsection{Data Collection Workflow}
Focus on the workflow for data we DID use, not everything we did

- gymkhana with stmpc, access to centerline

- talk episodes, forward/reverse, maps, speeds, binned lidar, curvature based lookahead, etc

- talk about periodic noise resulting from misaligned controller/sim timesteps

- logging package for both mpc and lstm captures state, and compute/latency metrics

\subsection{Model Training Workflow}

As data needs were defined and data was collected, an accompanying model training workflow was defined to reduce delays. Several training assumptions were defined to reduce the design space and focus on comparisons on model size and map performance in a narrow band. LSTMs were chosen as their internal logic gates that allowed for automatic rejection of poor or deformed inputs, such as noise in LiDAR scans. The general architecture of the LSTM model is shown in Figure \ref{model_arch}.

\begin{figure}[htbp]
    \centering
    \includegraphics[width=\linewidth]{assets/lstm_diagram.png}
    \caption{General ApproxiMPC Model Architecture}
    \label{model_arch}
\end{figure}

The main training assumptions for this study were that only LSTM models would be considered and that the input sequence length would be 100 steps and the output would be 1 timestep. This made training very consistent and reduced complexity within the scope of the project. 

The rest of the parameters were alterable including the LSTM hidden dimension, the number of LSTM blocks, and MLP Parameters. 

The standardized input features for the two models considered in this study are listed in Table \ref{tab:feature_comparison}.

\begin{table}[htbp]
    \renewcommand{\arraystretch}{1.3} % Adds standard IEEE vertical padding
    \caption{Comparison of Input Features Between Model V1 and Model V2}
    \label{tab:feature_comparison}
    \centering
    \begin{tabular}{l c c}
        \hline\hline
        \textbf{Attribute / Feature} & \textbf{Model V1} & \textbf{Model V2} \\
        \hline
        \textbf{Dataset} & Dataset 1 & Dataset 2 \\
        \textbf{Total Input Features} & 63 & 68 \\
        \textbf{Past Steering Angle} & Yes & Yes \\
        \textbf{Past Velocity Values} & Yes & Yes \\
        \textbf{Lidar Batches$^{\mathrm{a}}$} & Yes & Yes \\
        \textbf{Curvature Lookahead$^{\mathrm{b}}$} & No & Yes \\
        \hline
        \multicolumn{3}{l}{$^{\mathrm{a}}$Condensed to 60 min input buckets to reduce parameters.} \\
        \multicolumn{3}{l}{$^{\mathrm{b}}$Arc-length distance: 1, 2, 3, 5, 8 meters ahead.} \\
        \hline\hline
    \end{tabular}
\end{table}

The two models used in this study were iteratively trained to refine performance versus the MPC baseline. 

\textbf{Model V1:} Initial mapless model, used to validate data collection and training assumptions along with identifying early issues

\textbf{Model V2:} Improved Curvature model, used a larger dataset and map curvature inputs to improve performance for validation and comparison to MPC.

For training LSTMs and other sequence models, there are two possible training methods, teacher-forcing or autoregression. The teacher force passes the correct values at each subsequent time step during training, while the autoregression passes the predictions as input to the next time step \cite{davies2024teacherforcing}.

Due to the high dimensional inputs and low dimensional outputs of the LSTM model, teacher-forcing is the only realistic option for training. This simplifies training, but can make generalization more difficult when the model encounters out of distribution data. 

Talk deployment,
After training, 
convert to ONNX
Plug ONNX into a ROS package for inference in SIM.
Future implementations could utilize TensorRT for reduced latency and memory requirements. 

\subsection{Evaluation Methodology-Reggie}
Focus on the workflow for data we DID use, not everything we did

- discuss porting race performance to a "standard kmpc module with similar speed control to expert teacher"

- talk the 5 map comparison (why we used some from train AND test, simple to challenging, etc)

- talk the two map comparison for latency

- results from latency/compute analysis


\subsection{Results}

\subsubsection{Model V1 Results}

The initial V1 model was used to validate the training and testing assumptions with a mapless method. Another important aspect of Model V1 was determining acceptable model sizes for deployment.


- Chart of model values
- Size, num params and R2 values with colors 


These models were trained in Dataset V1 using teacher forcing as seen in Figure X.

- Training chart examples.

Discuss results from training and mention perturbation? (depends on data section)

After training was completed, an ROS package was made to test the performance of the Model V1 architectures on a subset of maps shown in Table X.

- Table of completed maps. 

To identify why Model V1 had difficulty completing full laps, the racelines must be analyzed. 


- Basic set of images of racelines from Model V1, if annoying just talk about it.


These raceline illustrate the limitations of Model V1: There are two main reasons for the failures of the V1 model, limits to the training data and limited knowledge of future track features. 

The limited perturbations and small dataset of Model V1 allowed rapid training of the models, but limited the distribution of data to which the LSTM was trained. This made the model poor in areas out of distribution when the model is off center. This causes failures on straightaways and turns when the model is unable to recognize it is off-center and react to return to the centerline. 

The other issue was limited knowledge of future track features. Due to Model V1 utilizing only the bagged LiDAR scans and pose information, it has limited knowledge of future conditions, like a hairpin coming up. This makes V1 react slowly or not at all to sharp turns.

\subsubsection{Model V2 Results}

These lessons from Model V1 were applied to the Dataset and Model V2 to improve performance versus the MPC baseline. The main changes to V2 were the use of a 10 times larger dataset with larger perturbations and the inclusion of lookahead sequence values as input to the model.

These values are the curvature of the centerline at 1, 2, 3, 5, and 8 m above the car. These inputs allow the LSTM to learn associated features such as hairpins or tight corners without needing to pass the full centerline information or directly use LiDAR scans. 

The results from Model V1 also suggested that models in the range of 125,000 parameters (roughly a hidden dimension of 128) had stronger perception. Due to this, the V2 models were set in a smaller band surrounding this value to further fine-tune the model size. 

- Chart of tested models per architecture. R2 values?

The training results of Model V2 in Figure X demonstrate improved model fitting, with final loss values almost half of Model V1. After training was complete, the models were tested on the same set of maps in ROS. The initial results of the tests can be seen in Figure X.

- Table of model successes

Once Model V2 was confirmed to pass 

- Lap times



- raceline comparisons 


- Latency comparisons 




- quantitative results

    - euclidean distance deviation from centerline, lap time, avg velocity
    
    - computation / memory / latency deltas

\section{Conclusion}


\section{Discussion and Future Work-Reggie}

- discuss limitations of using gym-khana
    - pros: fast iteration for data gen, configurable
    
    - cons: acados solver, not ROS compatible, stmpc technically
    
- no hardware validation

    - limitation of tool utilized
    
    - latency metrics attempt to compensate
    
- racing performance comparison not directly comparable

- two mpc models for performance vs compute

future work

- use kmpc bicycle model from getgo

- use gym khana for data gen, but add support for gym bridge

- run on hardware

- add speed control to raw data

- more work on GRU and RNNs

- lots I'm forgetting now I'm sure

- ADD GYM KHANA AS REFERENCE


\bibliographystyle{ieeetr}
\bibliography{references}

\end{document}
